\documentclass[oneside,openany,phd]{TMU-Thesis}


% ==================== structure/preamble.tex ====================
% مشخصات پایان‌نامه را در فایلهای faTitle و enTitle وارد نمایید.
% در اینجا برای یکپارچگی، همینجا وارد می‌کنیم

% مشخصات فارسی
\university{تربیت مدرس}
\faculty{مهندسی برق و کامپیوتر}
\department{گروه هوش مصنوعی و رباتیکز}
\subject{مهندسی کامپیوتر}
\field{مهندسی نرم‌افزار}
\credits{---} % در رساله دکتری معمولا به صورت واحدی ذکر نمی‌شود یا فیلد خالی است
\title{یادگیری تقویتی عمیق با شکل‌دهی تصادفی پاداش برای مدیریت عدم قطعیت در سیستم‌های خودوفقی}
\firstsupervisor{دکتر ...}
% \secondsupervisor{استاد راهنمای دوم}
\firstadvisor{دکتر ...}
% \secondadvisor{استاد مشاور دوم}
\name{علیرضا}
\surname{کاویانی‌فر}
\studentID{...}
\thesisdate{بهمن ۱۴۰۴}

% مشخصات انگلیسی
\latinuniversity{Tarbiat Modares University}
\latinfaculty{Faculty of Electrical and Computer Engineering}
\latinsubject{Computer Engineering}
\latinfield{Software Engineering}
\latintitle{Randomized Reward-Shaped Deep Reinforcement Learning for Managing Uncertainty in Self-Adaptive Systems}
\firstlatinsupervisor{Dr. ...}
\firstlatinadvisor{Dr. ...}
\latinname{Alireza}
\latinsurname{Kavianifar}
\latinthesisdate{February 2026}
\latinkeywords{Self-Adaptive Systems, Deep Reinforcement Learning, Randomized Reward Shaping, Uncertainty Management, Formal Verification, UPPAAL-SMC}

% فراخوانی بسته‌ها و دستورات استاندارد
% در این فایل، دستورها و تنظیمات مورد نیاز، آورده شده است.
%-------------------------------------------------------------------------------------------------------------------

% در ورژن جدید زی‌پرشین برای تایپ متن‌های ریاضی، این سه بسته، حتماً باید فراخوانی شود
\usepackage{amsthm,amssymb,amsmath}
% بسته‌ای برای تنطیم حاشیه‌های بالا، پایین، چپ و راست صفحه
\usepackage[top=30mm, bottom=25mm, left=25mm, right=35mm]{geometry}
% بسته‌‌ای برای ظاهر شدن شکل‌ها و تصاویر متن
\usepackage{graphicx}
% بسته‌ای برای رسم کادر
\usepackage{framed} 
% بسته‌‌ای برای چاپ شدن خودکار تعداد صفحات در صفحه «معرفی پایان‌نامه»
\usepackage{lastpage}
% بسته‌ و دستوراتی برای ایجاد لینک‌های رنگی با امکان جهش
% \usepackage[pagebackref=false,colorlinks,linkcolor=black,citecolor=black]{hyperref}
% چنانچه قصد پرینت گرفتن نوشته خود را دارید، خط بالا را غیرفعال و  از دستور زیر استفاده کنید چون در صورت استفاده از دستور زیر‌‌، 
% لینک‌ها به رنگ سیاه ظاهر خواهند شد که برای پرینت گرفتن، مناسب‌تر است
% \usepackage[pagebackref=false]{hyperref}
% بسته‌ لازم برای تنظیم سربرگ‌ها
\usepackage{fancyhdr}
\setlength{\headheight}{16pt}
%
\usepackage{setspace}
\usepackage{algorithm}
\usepackage{algorithmic}
\usepackage{subfigure}
\usepackage{titletoc}


% بسته‌ای برای ظاهر شدن «مراجع» و «نمایه» در فهرست مطالب
\usepackage[nottoc]{tocbibind}
% دستورات مربوط به ایجاد نمایه
\usepackage{makeidx}
\makeindex
% تعریف شیم‌های مورد نیاز برای نسخه‌های جدید بسته bidi در برخی محیط‌ها
% Removal of redundant and potentially broken definitions
\makeatletter
% شیم‌های سازگاری برای نسخه‌های جدید لاتک و بسته‌های bidi/xepersian
\providecommand{\IfClassLoadedT}[2]{\@ifclassloaded{#1}{#2}{}}
\providecommand{\IfClassLoadedF}[2]{\@ifclassloaded{#1}{}{#2}}
\providecommand{\IfClassLoadedTF}[3]{\@ifclassloaded{#1}{#2}{#3}}
\providecommand{\IfPackageLoadedT}[2]{\@ifpackageloaded{#1}{#2}{}}
\providecommand{\IfPackageLoadedF}[2]{\@ifpackageloaded{#1}{}{#2}}
\providecommand{\IfPackageLoadedTF}[3]{\@ifpackageloaded{#1}{#2}{#3}}
\providecommand{\@iffileloaded}[3]{\expandafter\ifx\csname ver@#1\endcsname\relax#3\else#2\fi}
\providecommand{\IfFileLoadedT}[2]{\@iffileloaded{#1}{#2}{}}
\providecommand{\IfFileLoadedF}[2]{\@iffileloaded{#1}{}{#2}}
\providecommand{\IfFileLoadedTF}[3]{\@iffileloaded{#1}{#2}{#3}}
\makeatother

%%%%%%%%%%%%%%%%%%%%%%%%%%
% فراخوانی بسته زی‌پرشین و تعریف قلم فارسی و انگلیسی
\usepackage[pagebackref=false,colorlinks,linkcolor=black,citecolor=black]{hyperref}
\usepackage{notoccite}

\ExplSyntaxOn
\cs_if_exist:NTF \g__hyp_linknestlevel_int { } { \int_new:N \g__hyp_linknestlevel_int }
\cs_if_exist:NTF \l__hyp_annot_GoTo_bool { } { \bool_new:N \l__hyp_annot_GoTo_bool }
\cs_if_exist:NTF \l__hyp_annot_URI_bool { } { \bool_new:N \l__hyp_annot_URI_bool }
\cs_if_exist:NTF \l__hyp_href_url_ismap_bool { } { \bool_new:N \l__hyp_href_url_ismap_bool }
\cs_if_exist:NTF \l__hyp_text_enc_uri_print_tl { } { \tl_new:N \l__hyp_text_enc_uri_print_tl }
\cs_if_exist:NTF \l__hyp_uri_tmpa_tl { } { \tl_new:N \l__hyp_uri_tmpa_tl }

% Defining missing l3pdf/pdfdict commands if not present
\cs_if_exist:NTF \pdfdict_put:nno { } { \cs_new:Npn \pdfdict_put:nno #1#2#3 { } }
\cs_if_exist:NTF \pdfdict_put:nnn { } { \cs_new:Npn \pdfdict_put:nnn #1#2#3 { } }
\cs_if_exist:NTF \pdfdict_use:n { } { \cs_new:Npn \pdfdict_use:n #1 { } }
\cs_if_exist:NTF \pdfannot_dict_put:nne { } { \cs_new:Npn \pdfannot_dict_put:nne #1#2#3 { } }
\cs_if_exist:NTF \pdfannot_link:nen { } { \cs_new:Npn \pdfannot_link:nen #1#2#3 { } }
\cs_if_exist:NTF \__hyp_check_link_nesting:TF { } { \cs_new:Npn \__hyp_check_link_nesting:TF #1#2 { #1 } }
\cs_if_exist:NTF \__hyp_text_pdfstring:eoN { } { \cs_new:Npn \__hyp_text_pdfstring:eoN #1#2#3 { } }
\ExplSyntaxOff

% تعریف دستورات گمشده برای جلوگیری از خطای renewcommand در بسته‌های داخلی
\makeatletter
\providecommand{\foottextfont}{}
\providecommand{\LTRfoottextfont}{}
\providecommand{\RTLfoottextfont}{}
\providecommand{\abstractname}{}
\providecommand{\latinabstract}{}
\providecommand{\AutoPersianDigits}{}
\makeatother

\usepackage{xepersian}
\settextfont[Path=font/,Extension=.ttf,Scale=1,
    BoldFont=XB NiloofarBd,
    ItalicFont=XB NiloofarIt,
    BoldItalicFont=XB NiloofarBdIt
]{XB Niloofar}
\setlatintextfont[Scale=0.9]{Times New Roman}

%%%%%%%%%%%%%%%%%%%%%%%%%%
% چنانچه می‌خواهید اعداد در فرمول‌ها، انگلیسی باشد، خط زیر را غیرفعال کنید
\setdigitfont[Path=font/,Extension=.ttf,Scale=1,
    BoldFont=XB NiloofarBd,
    ItalicFont=XB NiloofarIt,
    BoldItalicFont=XB NiloofarBdIt
]{XB Niloofar}
%%%%%%%%%%%%%%%%%%%%%%%%%%
% تعریف قلم‌های فارسی و انگلیسی اضافی برای استفاده در بعضی از قسمت‌های متن
\defpersianfont\titlefont[Path=font/,Extension=.ttf,Scale=1]{XB Titre}
% \defpersianfont\iranic[Scale=1.1]{XB Zar Oblique}%Italic}%
% \defpersianfont\nastaliq[Scale=1.2]{IranNastaliq}

%%%%%%%%%%%%%%%%%%%%%%%%%%
% دستوری برای حذف کلمه «چکیده»
\renewcommand{\abstractname}{}
% دستوری برای حذف کلمه «abstract»
%\renewcommand{\latinabstract}{}
% دستوری برای تغییر نام کلمه «اثبات» به «برهان»
\renewcommand\proofname{\textbf{برهان}}
% دستوری برای تغییر نام کلمه «کتاب‌نامه» به «مراجع»
\renewcommand{\bibname}{مراجع}
% دستوری برای تعریف واژه‌نامه انگلیسی به فارسی
\newcommand\persiangloss[2]{#1\dotfill\lr{#2}\\}
% دستوری برای تعریف واژه‌نامه فارسی به انگلیسی 
\newcommand\englishgloss[2]{#2\dotfill\lr{#1}\\}
% تعریف دستور جدید «\پ» برای خلاصه‌نویسی جهت نوشتن عبارت «پروژه/پایان‌نامه/رساله»
\newcommand{\پ}{پروژه/پایان‌نامه/رساله }

%\newcommand\BackSlash{\char`\\}

%%%%%%%%%%%%%%%%%%%%%%%%%%
\SepMark{-}

% تعریف و نحوه ظاهر شدن عنوان قضیه‌ها، تعریف‌ها، مثال‌ها و ...
\theoremstyle{definition}
\newtheorem{definition}{تعریف}[section]
\theoremstyle{plain}
\newtheorem{theorem}[definition]{قضیه}
\newtheorem{lemma}[definition]{لم}
\newtheorem{proposition}[definition]{گزاره}
\newtheorem{corollary}[definition]{نتیجه}
\newtheorem{remark}[definition]{ملاحظه}
\theoremstyle{definition}
\newtheorem{example}[definition]{مثال}

%\renewcommand{\theequation}{\thechapter-\arabic{equation}}
%\def\bibname{مراجع}
\numberwithin{algorithm}{chapter}
\def\listalgorithmname{فهرست الگوریتم‌ها}
\def\listfigurename{فهرست تصاویر}
\def\listtablename{فهرست جداول}

%%%%%%%%%%%%%%%%%%%%%%%%%%%%
% دستورهایی برای سفارشی کردن سربرگ صفحات
% \newcommand{\SetHeader}{
% \csname@twosidetrue\endcsname
% \pagestyle{fancy}
% \fancyhf{} 
% \fancyhead[OL,EL]{\thepage}
% \fancyhead[OR]{\small\rightmark}
% \fancyhead[ER]{\small\leftmark}
% \renewcommand{\chaptermark}[1]{%
% \markboth{\thechapter-\ #1}{}}
% }
%%%%%%%%%%%%5
%\def\MATtextbaseline{1.5}
%\renewcommand{\baselinestretch}{\MATtextbaseline}
\doublespacing
%%%%%%%%%%%%%%%%%%%%%%%%%%%%%
% دستوراتی برای اضافه کردن کلمه «فصل» در فهرست مطالب با استفاده از titletoc
\titlecontents{chapter}
  [5em]
  {\addvspace{1em}\bfseries}
  {\contentslabel[\chaptername~\thecontentslabel:]{5em}}
  {\hspace*{-5em}}
  {\titlerule*[0.5pc]{.}\contentspage}

\titlecontents{section}
  [5em]
  {}
  {\contentslabel{2em}}
  {\hspace*{-2em}}
  {\titlerule*[0.5pc]{.}\contentspage}

\titlecontents{subsection}
  [8em]
  {}
  {\contentslabel{3em}}
  {\hspace*{-3em}}
  {\titlerule*[0.5pc]{.}\contentspage}

\titlecontents{subsubsection}
  [12em]
  {}
  {\contentslabel{4em}}
  {\hspace*{-4em}}
  {\titlerule*[0.5pc]{.}\contentspage}

\makeatletter
{\def\thebibliography#1{\chapter*{\refname\@mkboth
   {\uppercase{\refname}}{\uppercase{\refname}}}\list
   {[\arabic{enumi}]}{\settowidth\labelwidth{[#1]}
   \rightmargin\labelwidth
   \advance\rightmargin\labelsep
   \advance\rightmargin\bibindent
   \itemindent -\bibindent
   \listparindent \itemindent
   \parsep \z@
   \usecounter{enumi}}
   \def\newblock{}
   \sloppy
   \sfcode`\.=1000\relax}}
\makeatother


% تنظیمات اختصاصی برای فهرست مطالب (TOC)
\setcounter{secnumdepth}{3}
\setcounter{tocdepth}{3}

\raggedbottom

% تنظیمات نقطه‌چین برای فهرست مطالب
\usepackage{tocloft}
\providecommand{\cftchapdotfill}{\cftdotfill{\cftdotsep}}
\providecommand{\cftsecdotfill}{\cftdotfill{\cftdotsep}}
\providecommand{\cftsubsecdotfill}{\cftdotfill{\cftdotsep}}


% تنظیمات سربرگ برای رفع مشکل نمایش حروف فارسی (جعبه‌های خالی)
\usepackage{fancyhdr}
\pagestyle{fancy}
\fancyhf{}
\fancyhead[L]{\thepage}
\fancyhead[R]{\small\leftmark}
\renewcommand{\chaptermark}[1]{\markboth{\chaptername\ \thechapter.\ #1}{}}
\renewcommand{\sectionmark}[1]{\markright{\thesection.\ #1}}
% حذف حالت ایتالیک/کج که باعث به هم ریختگی فونت فارسی در سربرگ می‌شود
\let\oldheadrule\headrule
\renewcommand{\headrule}{\oldheadrule}
\fancyhead[R]{\small\nouppercase{\leftmark}}



\begin{document}
\AutoPersianDigits
\setstretch{1.5}

\setcounter{chapter}{2}

\chapter{مرور پیشینه پژوهش و کارهای مرتبط}

\section{مقدمه و روش‌شناسی مرور}
\label{sec:ch3_intro}

هدف بنیادین این فصل، ارائه یک بررسی نظام‌مند، انتقادی و عمیق از ادبیات علمی و رویکردهای موجود در زمینه مدیریت عدم قطعیت و تصمیم‌گیری زمان اجرا در سیستم‌های خودوفق (\lr{Self-Adaptive Systems - SAS}) است. در طراحی مهندسی سیستم‌های خودوفق مدرن، اتخاذ تصمیمات وفقی بهینه در مواجهه با شرایط پویای محیطی و فضاهای تصمیم‌گیری بسیار بزرگ، یکی از چالش‌های بنیادین به شمار می‌رود. کارهای پیشین هر یک از زاویه‌ای خاص به این مسئله نگریسته‌اند؛ با این حال، به منظور درک دقیق سهم علمی این رساله و تبیین دقیق ضرورت ارائه الگوریتم \lr{RS-DRL}، تحلیل موشکافانه و مقایسه‌ای این رویکردها امری حیاتی است.

بدین منظور، در این فصل روش‌شناسی مروری ساختاریافته‌ای اتخاذ گردیده است. فرآیند شناسایی و گزینش مقالات از طریق جستجوی جامع در پایگاه‌های داده معتبر بین‌المللی از جمله \lr{IEEE Xplore}، \lr{ACM Digital Library}، \lr{Scopus} و \lr{Google Scholar} صورت پذیرفته است. بازه زمانی جستجو متمرکز بر پژوهش‌های برجسته دو دهه اخیر بوده و کلمات کلیدی نظیر «سیستم‌های خودوفق»، «یادگیری تقویتی آنلاین»، «شکل‌دهی پاداش»، «فضای وفقی بزرگ» و «تأیید صوری زمان اجرا» به عنوان معیارهای اصلی فیلترینگ استفاده شده‌اند. معیارهای ورود مقالات به چرخه بررسی انتقادی شامل تمرکز بر روی مدیریت فضاهای وفقی گسسته در زمان اجرا و توانایی سیستم در تصمیم‌گیری آنلاین تحت شرایط عدم قطعیت بوده است.

جهت ارزیابی نظام‌مند و مقایسه‌ی عادلانه پیشینه پژوهش، یک چارچوب ارزیابی پنج‌بعدی بر پایه معیارهای مهندسی نرم‌افزار خودوفق \cite{cheng2009software, weyns2012formal} توسعه داده شده است که تمامی آثار بررسی‌شده بر اساس این شاخص‌ها مورد نقد قرار می‌گیرند:
\begin{enumerate}
    \item \textbf{ابعاد فضای وفقی (\lr{Adaptation Space Size}):} توانایی متدولوژی پیشنهادی در مقیاس‌پذیری و مدیریت فضاهای تصمیم‌گیری گسسته و بزرگ (فراتر از ۱۰۰۰ گزینه وفقی).
    \item \textbf{مدیریت پویایی و عدم‌قطعیت زمان‌اجرا (\lr{Handling of Non-Stationarity}):} پایداری عملکرد سیستم در مواجهه با رانش مفهومی (\lr{Concept Drift}) و تغییرات پیش‌بینی‌نشده محیطی بدون افت کیفیت خدمات.
    \item \textbf{عدم وابستگی به دانش دامنه خاص (\lr{Need for Domain Knowledge}):} میزان نیاز به تخصص انسانی زمان توسعه یا طراحی اهداف صریح و پیچیده برای عامل یادگیرنده.
    \item \textbf{مقیاس‌پذیری و کارایی محاسباتی (\lr{Real-Time Scalability}):} توانایی تصمیم‌گیری آنلاین در محدوده میلی‌ثانیه به منظور رعایت قیود پایداری زمانی سیستم.
    \item \textbf{قابلیت خروج از بهینه‌های محلی (\lr{Ability to Escape Local Optima}):} کارایی مکانیزم‌های کاوش در فرار از تله‌های بهینگی محلی در طول یادگیری آنلاین.
\end{enumerate}\footnote{این پنج معیار ارزیابی برای اولین بار در این رساله تدوین شده‌اند. جدول ۱-۳ در انتهای همین فصل (صفحه ۴۰) مقایسه تطبیقی جامع روش پیشنهادی \lr{RS-DRL} را با کارهای پیشین بر اساس این پنج بعد ارائه می‌دهد.}

در ادامه‌ی این فصل، ابتدا مبانی نظری عدم قطعیت و تصمیم‌گیری در سیستم‌های خودوفق تشریح شده، سپس سه پارادایم اصلی تصمیم‌گیری (مبتنی بر یادگیری ماشین سنتی، مبتنی بر جستجو و مبتنی بر یادگیری تقویتی عمیق) به صورت موشکافانه تحلیل می‌گردند. در نهایت، با بررسی روش‌های پیشرفته بازپخش تجربه و شکل‌دهی پاداش و نقش روش‌های صوری، خلاءهای پژوهشی موجود استخراج و جایگاه الگوریتم پیشنهادی \lr{RS-DRL} تبیین می‌گردد.

\section{مبانی نظری: عدم‌قطعیت و تصمیم‌گیری در سیستم‌های خودوفق}
\label{sec:ch3_foundations}

سیستم‌های خودوفق به منظور حفظ اهداف کیفی خود (نظیر قابلیت اطمینان، مصرف انرژی و تأخیر)، نیازمند پایش مستمر محیط فیزیکی و تصمیم‌گیری بهینه در زمان اجرا هستند \cite{kephart2003vision, cheng2009software}. با این حال، تصمیم‌گیری در این سیستم‌ها تحت تأثیر شدید منابع گوناگون عدم قطعیت قرار دارد. عدم قطعیت در ادبیات مهندسی نرم‌افزار خودوفق \cite{cheng2009software, esfahani2013uncertainty} به طور کلی به سه دسته اصلی تقسیم می‌شود: عدم قطعیت محیطی (مانند نوسانات تداخل سیگنال در شبکه‌های اینترنت اشیا)، عدم قطعیت سیستمی (نظیر خرابی‌های ناگهانی سخت‌افزار یا تغییرات زیرساختی) و عدم قطعیت اهداف (تغییر در اولویت‌های کیفی کاربران در طول زمان).

برای مدل‌سازی فرآیند تصمیم‌گیری در مواجهه با این عدم قطعیت‌ها، مهندسان سیستم‌های خودوفق از ابزارهای ریاضیاتی متعددی بهره می‌برند:
\begin{itemize}
    \item \textbf{فرآیند تصمیم‌گیری مارکوف (\lr{MDP}):} استانداردترین روش مدل‌سازی ریاضی تصمیم‌گیری پیاپی است که در آن فرض می‌شود وضعیت بعدی سیستم تنها به وضعیت فعلی و اقدام اتخاذ شده وابسته است \cite{puterman1994markov}. در سیستم‌های خودوفق، وضعیت‌ها نشان‌دهنده شرایط محیطی و پیکربندی‌های سیستم، و اقدامات نشان‌دهنده گزینه‌های وفقی هستند.
    \item \textbf{فرآیند تصمیم‌گیری مارکوف نیمه‌مشاهده‌پذیر (\lr{POMDP}):} در شرایطی استفاده می‌شود که عامل یادگیرنده دسترسی کامل به وضعیت واقعی محیط ندارد و تنها مشاهدات نویزی در اختیار دارد \cite{kaelbling1998planning}. اگرچه \lr{POMDP} واقع‌گرایانه‌تر است، اما حل آن سربار محاسباتی بسیار بالایی دارد که مانع تصمیم‌گیری آنلاین بلادرنگ می‌شود.
    \item \textbf{جستجوی تصادفی زمان‌اجرا:} رویکردی غیرمدل‌محور است که سیستم با نمونه‌برداری زمان‌اجرا از فضای تغییرات، تلاش می‌کند گزینه‌ای را بیابد که اهداف سیستم را ارضا نماید.
\end{itemize}

یکی از چالش‌های بنیادین ریاضیاتی در تصمیم‌گیری آنلاین سیستم‌های خودوفق، پدیده \textbf{تله بهینه‌های محلی (\lr{Local Optima Trap})} است \cite{du2023scalability}. در فضاهای وفقی بزرگ و گسسته، تابع مطلوبیت یا پاداش سیستم معمولاً غیرخطی، چندهدفه و دارای تپه‌های متعدد است. هنگامی که یک الگوریتم تصمیم‌گیری آنلاین یا عامل یادگیرنده تلاش می‌کند تا از طریق کاوش محلی (مانند شیب صعودی یا استراتژی‌های حریصانه) پیکربندی بهینه را بیابد، به شدت مستعد گرفتار شدن در نقاطی است که عملکرد بهتری نسبت به همسایگان بلاواسطه خود دارند، اما با بهینه سراسری (\lr{Global Optima}) فاصله زیادی دارند.

این مشکل در شبکه‌های پویا نظیر \lr{DeltaIoT} تشدید می‌شود \cite{gerostathopoulos2019deltaiot}؛ زیرا در این شبکه‌ها، تغییرات تداخل محیطی باعث تغییر مداوم قله‌های بهینگی می‌شود. اگر عامل تصمیم‌گیرنده نتواند با ترکیب بهینه‌ی کاوش (\lr{Exploration}) و بهره‌برداری (\lr{Exploitation})، به طور مداوم و پویا فضای تصمیم‌گیری را جستجو کند، در یک سیاست تطبیقی منجمد شده و با تغییر شرایط محیطی دچار افت شدید کیفیت خدمات می‌گردد. در بخش‌های بعدی، نحوه مواجهه رویکردهای کلاسیک با این چالش ریاضیاتی را به تفصیل نقد خواهیم نمود.

\section{رویکردهای اصلی مدیریت فضای وفقی در زمان اجرا}
\label{sec:ch3_approaches}

به منظور تحلیل و مقایسه‌ی پیشینه پژوهش در زمینه مدیریت فضاهای وفقی بزرگ، آثار علمی موجود را می‌توان به سه پارادایم اصلی تقسیم‌بندی نمود: رویکردهای مبتنی بر یادگیری ماشین سنتی برای هرس فضا، روش‌های مبتنی بر جستجو و بهینه‌سازی آنلاین، و مدل‌های مبتنی بر یادگیری تقویتی برای تصمیم‌گیری زمان اجرا. در ادامه، هر یک از این رویکردها را بر اساس چارچوب ارزیابی پنج‌بعدی معرفی‌شده نقد علمی می‌نماییم.

\subsection{هرس فضای وفقی مبتنی بر یادگیری ماشین سنتی}
\label{sec:ch3_ml_pruning}

یکی از ایده‌های نخستین برای مقابله با انفجار ترکیبی گزینه‌های وفقی در سیستم‌های خودوفق بزرگ، هرس پیشاپیش فضای وفقی پیش از ارسال به مؤلفه‌ی تصمیم‌گیر (یا برنامه‌ریز) است. در این دسته از رویکردها، مهندسان سیستم تلاش می‌کنند تا با استفاده از مدل‌های پیش‌بینی‌کننده یادگیری ماشین سنتی، گزینه‌های نامطلوب را فیلتر نمایند.

پژوهش‌های برجسته‌ی کوئین (\lr{Quin}) و همکارانش \cite{quin2019efficient, quin2022reducing} از رویکردهای پیشگام در این زمینه است. آن‌ها با به‌کارگیری تکنیک‌های طبقه‌بندی نظارت‌شده (مانند درخت تصمیم و جنگل تصادفی) نشان دادند که می‌توان با آموزش یک مدل روی داده‌های تاریخی، پیکربندی‌هایی را که شانس بسیار پایینی برای ارضای کیفیت خدمات (\lr{QoS}) دارند، شناسایی و هرس کرد. به طور مشابه، کاچی و بواناکا \cite{kachi2023hybrid} متدولوژی \lr{ML4EAS} (\lr{Machine Learning for Evaluating Adaptation Space}) را توسعه دادند که از ترکیب الگوریتم‌های طبقه‌بندی و رگرسیون برای پیش‌بینی و غربالگری گزینه‌های ناکارآمد در سیستم‌های توزیع‌شده بهره می‌برد. این روش‌ها قادرند ابعاد فضای جستجو را در سناریوهای ایستا تا ۶۰ درصد کاهش دهند \cite{kachi2023hybrid}.

\textbf{نقد و بررسی انتقادی:}
علیرغم کارایی این روش‌ها در کاهش اندازه محیط تصمیم‌گیری زمان اجرا، دو محدودیت بنیادین در این متدولوژی‌ها وجود دارد.
از یک سو، این روش‌ها با \textbf{وابستگی شدید به داده‌های برچسب‌دار آفلاین} مواجه هستند؛ چرا که فرآیند آموزش این طبقه‌بندها مستلزم جمع‌آوری مجموعه‌داده‌های بسیار حجیم و متنوع در زمان توسعه (\lr{Design-time}) است و در سیستم‌های پیچیده معاصر، شبیه‌سازی تمامی سناریوهای نویزی و محیطی در زمان طراحی ناممکن است \cite{quin2022reducing, gheibi2024dealing}.
از سوی دیگر، \textbf{پدیده رانش مفهومی (\lr{Concept Drift})} چالش مهم دیگری است؛ غیبی و وینز \cite{gheibi2024dealing} در تحلیل‌های خود تصریح کرده‌اند که در محیط‌های بسیار پویا (مانند اینترنت اشیا شهری با تداخلات متغیر)، الگوهای نویز و ویژگی‌های سیستم دستخوش تغییر مستمر زمان اجرا می‌شوند. با رخ دادن رانش مفهومی، فرضیات اولیه طبقه‌بندهای یادگیری ماشین نامعتبر گشته و دقت مدل در هرس گزینه‌ها به شدت افت می‌کند. برای جبران این افت، سیستم ناگزیر به آموزش مجدد (\lr{Retraining}) مکرر و بسیار پرهزینه مدل در زمان اجراست، که این فرآیند علاوه بر ایجاد سربار محاسباتی زیاد، سیستم را در طول فرآیند آموزش در وضعیت ناپایدار و غیرایمن قرار می‌دهد و در نهایت منجر به ایجاد یک \textbf{موازنه شکننده بین مقیاس‌پذیری و استواری} می‌گردد.

\subsection{روش‌های مبتنی بر جستجوی آنلاین و استراتژی‌های اکتشافی}
\label{sec:ch3_search_based}

پارادایم دوم، تصمیم‌گیری زمان اجرا را به عنوان یک مسئله‌ی بهینه‌سازی و جستجوی آنلاین مدل‌سازی می‌کند. در این روش‌ها، نیازی به داده‌های پیش‌آموزش دیده نیست؛ بلکه سیستم با وقوع هرگونه تغییر در محیط، الگوریتم‌های جستجو را برای یافتن بهترین پیکربندی در فضای وفقی اجرا می‌کند.

کوکر (\lr{Coker}) در معماری \lr{SASS} \cite{coker2015sass} از روش‌های جستجوی موضعی تصادفی نظیر الگوریتم صعود تپه (\lr{Hill Climbing}) استفاده کرد تا با بررسی همسایگی پیکربندی فعلی، به پیکربندی برتر دست یابد. اگرچه این روش محاسبات سبکی دارد، اما به شدت مستعد گرفتار شدن در اولین بهینه محلی است \cite{coker2015sass}. برای عبور از این چالش، پژوهش‌های کینر (\lr{Kinneer}) و همکارانش \cite{kinneer2018managing, kinneer2020building, kinneer2021information} بر به‌کارگیری برنامه‌نویسی ژنتیک (\lr{Genetic Programming}) و الگوریتم‌های تکاملی زمان اجرا متمرکز شد. آن‌ها نشان دادند که با تکامل جمعیت‌های گزینه‌های وفقی می‌توان به راه‌حل‌های بسیار بهتری دست یافت. علاوه بر این، چن (\lr{Chen}) در رویکرد \lr{FEMOSAA} \cite{chen2018femosaa} با استفاده از الگوریتم تکاملی چندهدفه (\lr{MOEA})، مصالحه‌ی پویای بین اهداف متناقض کیفی (مانند انرژی و تأخیر) را در زمان اجرا مدل‌سازی و جستجو نمود.

\textbf{نقد و بررسی انتقادی:}
اگرچه این رویکردها ویژگی تطبیق‌پذیری آنلاین بدون نیاز به آموزش را فراهم می‌سازند، اما نقطه ضعف بزرگ آن‌ها \textbf{سربار محاسباتی زمان اجرا و چالش پایداری زمانی} است:
اولین نقطه ضعف، \textbf{عدم مقیاس‌پذیری در ابعاد بزرگ} است؛ الگوریتم‌های تکاملی و چندهدفه برای همگرایی به جبهه بهینه پارتو (\lr{Pareto Front}) نیازمند ارزیابی صدها یا هزاران پیکربندی در هر گام تصمیم‌گیری هستند. هنگامی که ابعاد فضای وفقی بزرگ شده و فراتر از ۱۰۰۰ گزینه (نظیر فضای ۴۰۹۶ گزینه‌ای بستر \lr{DeltaIoTv2}) می‌رود، زمان محاسبات این الگوریتم‌ها به مقیاس ثانیه یا دقیقه افزایش می‌یابد \cite{coker2015sass, kinneer2021information}. دومین نقطه ضعف، \textbf{تأخیر غیرقابل‌تحمل زمان تصمیم‌گیری} است؛ در سیستم‌های حساس به زمان، این تأخیر طولانی برنامه‌ریز منجر به بروز رفتارهای مخرب و انحراف شدید از اهداف پایداری سیستم در طول بازه تصمیم‌گیری می‌گردد، که استقرار عملیاتی این روش‌ها را در محیط‌های زمان‌حقیقی غیرممکن می‌سازد \cite{coker2015sass, kinneer2021information}.

\subsection{یادگیری تقویتی زمان‌اجرا و کنترل خودمختار}
\label{sec:ch3_rl_adaptation}

یادگیری تقویتی (\lr{RL}) به عنوان یک پارادایم تصمیم‌گیری آنلاین پویا، به دلیل توانایی ذاتی عامل در یادگیری سیاست وفقی مطلوب از طریق تعامل مستقیم با محیط، توجه بسیاری را جلب نموده است. 

کارهای کلاسیک کیم و پارک \cite{kim2009reinforcement} از اولین تلاش‌ها در به‌کارگیری الگوریتم \lr{Q-Learning} جدول‌محور برای تطابق پویای نرم‌افزارها بود. در ادامه، متزگر و همکارانش \cite{metzger2022realizing} قابلیت یادگیری تقویتی را برای مدیریت تطابق در معماری‌های سرویس‌گرا تعمیم دادند. همچنین، کامارا و همکارانش \cite{camara2020quantitative} چارچوب \lr{RLMC} (\lr{Reinforcement Learning Model Checking}) را معرفی کردند که در آن از یادگیری تقویتی آنلاین برای هدایت سیستم‌های خودوفق استفاده شده و ایمنی کاوش را از طریق مدل‌چکینگ تضمین می‌کرد. با این حال، تمامی این کارهای اولیه به دلیل استفاده از روش‌های جدول‌محور، با چالش جدی «طاعون بعدنمایی» (\lr{Curse of Dimensionality}) در فضاهای حالت و عمل بزرگ مواجه بودند \cite{sutton2018reinforcement}.

با ظهور یادگیری تقویتی عمیق (\lr{DRL}) \cite{mnih2015human}، محققان تلاش کردند تا با استفاده از شبکه‌های عصبی عمیق به عنوان تقریب‌زننده‌های تابع ارزش (مانند معماری عمیق \lr{DQN})، فضاهای بزرگ سیستم‌های خودوفق را مدیریت کنند \cite{du2023scalability}.

\textbf{نقد و بررسی انتقادی:}
با وجود توانایی شبکه‌های عمیق در تقریب توابع در فضاهای بزرگ، پیاده‌سازی استانداردهای یادگیری تقویتی عمیق (نظیر \lr{DQN} با مکانیزم اکتشاف حریصانه‌ی سنتی \lr{$\epsilon$-greedy}) در سیستم‌های خودوفق با یک چالش تئوریک و عملی بزرگ مواجه است.
شاخص‌ترین نقطه ضعف این رویکردها، \textbf{تله عمیق بهینه‌های محلی (\lr{Local Optima Trap})} است؛ دو (\lr{Du}) و همکارانش \cite{du2023scalability} در تحلیل‌های خود نشان دادند که در فضاهای وفقی بزرگ و گسسته، پویایی‌های نویزی زمان اجرا باعث می‌شوند که سیاست عامل یادگیرنده به راحتی در بهینه‌های محلی متوقف شود. استراتژی کاوش تصادفی \lr{$\epsilon$-greedy} در مواجهه با ابعاد بسیار بزرگ (مانند ۴۰۹۶ گزینه وفقی)، درصد بسیار بالایی از کاوش خود را صرف گزینه‌های غیرمفید و مخرب می‌کند. این امر همگرایی عامل را به شدت کند کرده و مدل را وادار می‌سازد تا در اولین پیکربندی نیمه‌پایدار متوقف گشته و توانایی عبور و رسیدن به بهینه سراسری را نداشته باشد، مگر آنکه فرآیند آموزش مجدد مدل با مقادیر بهینه‌تر از ابتدا آغاز گردد که در زمان اجرا به شدت پرهزینه است.

\section{تکنیک‌های پیشرفته یادگیری تقویتی: کاوش، شکل‌دهی پاداش و نگاه به گذشته}
\label{sec:ch3_advanced_rl}

در سیستم‌های خودوفق معاصر، عبور از محدودیت‌های الگوریتم‌های یادگیری تقویتی سنتی مستلزم به‌کارگیری تکنیک‌های پیشرفته در سه حوزه کلیدی است: مدیریت فرآیند کاوش در فضاهای وسیع، استفاده از مکانیزم‌های بازپخش تجربه با نگاه به گذشته، و استراتژی‌های شکل‌دهی پاداش. در این بخش، مبانی نظری و ریاضیاتی هر یک از این حوزه‌ها را تبیین کرده و چالش‌های تعمیم آن‌ها به سیستم‌های خودوفق را به صورت انتقادی واکاوی می‌نماییم.

\subsection{چالش‌های استراتژی کاوش در فضاهای وفقی بزرگ}
\label{sec:ch3_exploration_challenges}

فرآیند کاوش (\lr{Exploration}) در یادگیری تقویتی، مکانیزمی است که به عامل اجازه می‌دهد تا رفتارهای جدید را به منظور کشف پاداش‌های بالاتر بیازماید. در فضاهای وفقی بزرگ و گسسته (مانند شبکه \lr{DeltaIoT} با ۴۰۹۶ پیکربندی ممکن)، طراحی استراتژی کاوش به یک چالش بسیار پیچیده‌ی ریاضی تبدیل می‌شود.

روش‌های کاوش سنتی نظیر \lr{$\epsilon$-greedy} یا کاوش بولتزمن (\lr{Boltzmann Exploration}) در فضاهای با ابعاد بالا به شدت ناکارآمد هستند \cite{du2023scalability, du2023alleviating}. در استراتژی \lr{$\epsilon$-greedy}، عامل با احتمال $\epsilon$ اقدام بعدی خود را به صورت کاملاً تصادفی انتخاب می‌کند. هنگامی که تعداد اقدامات ممکن $A = 4096$ باشد، احتمال انتخاب یک اقدام وفقی مفید در میان انبوهی از اقدامات غیربهینه یا مخرب، بسیار ناچیز است. در نتیجه، بخش عمده‌ای از گام‌های کاوش عامل صرف بررسی گزینه‌هایی می‌شود که کیفیت خدمات را به شدت کاهش داده و پایداری سیستم را به مخاطره می‌اندازند. این امر منجر به افت شدید «کارایی نمونه» (\lr{Sample Efficiency}) شده و سرعت همگرایی آنلاین را تا حد غیرقابل‌تحملی کاهش می‌دهد.

استفاده از روش‌های پیشرفته‌تر مانند حد بالای اطمینان (\lr{Upper Confidence Bound - UCB}) نیز در محیط‌های زمان اجرا با عدم قطعیت پویا با مشکل مواجه است؛ زیرا ارزیابی مستمر عدم قطعیتِ تک‌تکِ اقدامات در فضاهای بسیار بزرگ، سربار محاسباتی شدیدی را به کنترل‌کننده زمان اجرا تحمیل می‌کند \cite{sutton2018reinforcement}. از این رو، سیستم‌های خودوفق هوشمند نیازمند یک استراتژی کاوش هدایت‌شده هستند که بدون نیاز به نمونه‌برداری تصادفی کورکورانه، عامل را به سمت مناطق امیدبخش فضای وفقی هدایت نماید.

\subsection{ناکارآمدی مکانیزم بازپخش تجربه با نگاه به گذشته \lr{(HER)} در محیط‌های آستانه‌محور}
\label{sec:ch3_her_inefficiency}

مکانیزم بازپخش تجربه با نگاه به گذشته (\lr{Hindsight Experience Replay - HER}) که توسط آندریچویچ (\lr{Andrychowicz}) و همکارانش \cite{andrychowicz2017her} معرفی شد، یکی از برجسته‌ترین پیشرفت‌ها در زمینه حل مسئله‌ی پاداش‌های پراکنده (\lr{Sparse Rewards}) است. ایده کلیدی در \lr{HER}، برچسب‌گذاری مجدد هدف (\lr{Goal Relabeling}) برای تجربیات ناموفق است. به طور مشابه، رویکرد پاداش نرم‌تر (\lr{Soft HER}) توسعه‌یافته توسط هه (\lr{He}) و همکارانش \cite{he2020softher} تلاش نمود تا فرآیند برچسب‌گذاری مجدد را با استفاده از پاداش‌های ملایم‌تر سرعت بخشد.

با این وجود، بر اساس تحلیل تئوریک ارائه شده در این رساله، استدلال می‌کنیم که این مکانیزم‌ها دارای یک \textbf{تناقض و ناسازگاری بنیادین با پارادایم خودوفقی} هستند:
\begin{itemize}
    \item \textbf{طراحی هدف‌محور (\lr{Goal-Conditioned}) در مقابل ماهیت آستانه‌محور (\lr{Threshold-Based}) خودوفقی:} الگوریتم \lr{HER} به طور ذاتی برای محیط‌های هدف‌محور طراحی شده است؛ محیط‌هایی که در آن‌ها هدف عامل، رسیدن به یک مختصات یا وضعیت فیزیکی مشخص و صریح (مانند مکان سه بعدی بازوی رباتیک $g \in G$) است. در این سناریوها، اگر عامل نتواند به هدف $g$ برسد، \lr{HER} وضعیت نهایی تجربه شکست‌خورده ($s'$) را به عنوان هدف فرضی در نظر گرفته و پاداش مثبتی به آن اختصاص می‌دهد.
    \item \textbf{فقدان نقطه هدف صریح در سیستم‌های خودوفق:} اهداف سیستم‌های خودوفق به ندرت به صورت یک نقطه وضعیت واحد فرموله می‌شوند. در مقابل، این اهداف همواره \textbf{آستانه‌محور} هستند. برای مثال، هدف سیستم کنترلی شبکه \lr{DeltaIoT}، نگه‌داشتن میزان تأخیر در کمتر از ۱۰۰ میلی‌ثانیه و مصرف انرژی در کمتر از ۲ وات است. در این حالت، هدف یک نقطه مشخص نیست، بلکه زیرمجموعه وسیعی از فضا است که قیود آستانه‌ای را ارضا می‌سازد.
    \item \textbf{عدم امکان برچسب‌گذاری مجدد:} هنگامی که یک عامل در زمان اجرا شکست می‌خورد (مثلاً تأخیر شبکه به ۱۵۰ میلی‌ثانیه می‌رسد و آستانه نقض می‌شود)، هیچ مختصات یا وضعیت واحدی وجود ندارد که بتوان شکست را به عنوان یک «هدف جایگزین معتبر» برچسب‌گذاری مجدد کرد. نقض آستانه اساساً به معنای ناامنی و افت کیفیت خدمات است و تعریف یک «آستانه مجازی نقض‌شده» به عنوان هدف جدید در یادگیری با نگاه به گذشته، از نظر مهندسی و ریاضیاتی فاقد اعتبار بوده و ایمنی و ثبات کنترل‌کننده را به طور کامل نابود می‌سازد.
\end{itemize}

بنابراین، تکنیک‌های هدف‌محوری مانند \lr{HER} و \lr{Soft HER} عملاً قابلیت کاربرد در کنترل‌کننده‌های سیستم‌های خودوفق آستانه‌محور را دارا نمی‌باشند\footnote{با این حال، در سیستم‌های خودوفق آستانه‌محور، استفاده از توابع هدف صریح یا پتانسیل‌های ثابت با چالش ناسازگاری مواجه است. این ناسازگاری یکی از شکاف‌های تحقیقاتی اصلی این رساله را تشکیل می‌دهد. راهکار پیشنهادی \lr{RS-DRL} در فصل ۴ (بخش ۴-۶) با مکانیسم شکل‌دهی تصادفی پاداش (بدون نیاز به هدف صریح) به این چالش پاسخ می‌دهد. نتایج تجربی برتری آن نسبت به \lr{HER} در فصل ۵ (جداول \ref{tab:min_goals_v1} و \ref{tab:min_goals_v1_1}) تأیید شده است.}.

\subsection{مبانی نظری و ریاضیاتی شکل‌دهی پاداش}
\label{sec:ch3_reward_shaping_math}

شکل‌دهی پاداش (\lr{Reward Shaping}) رویکردی قدرتمند در نظریه یادگیری تقویتی برای تسریع همگرایی عامل از طریق افزودن یک سیگنال پاداش کمکی (\lr{Auxiliary Reward}) به پاداش اصلی محیط است. از نظر ریاضیاتی، پاداش شکل‌دهی‌شده زمان اجرا به صورت زیر تعریف می‌شود:
\begin{equation}
    R_{shaped}(s, a, s') = R_{env}(s, a, s') + F(s, a, s')
\end{equation}
که در آن $R_{env}$ پاداش حاصل از دینامیک محیط و $F(s, a, s')$ پاداش کمکی شکل‌دهی است.

نگ (\lr{Ng}) و همکارانش در پژوهش کلاسیک خود \cite{ng1999policy} اثبات ریاضیاتی ارائه‌دادند که نشان می‌دهد شکل‌دهی پاداش کورکورانه می‌تواند منجر به بروز پدیده «چرخه‌های بهینگی کاذب» (\lr{Suboptimal Loops}) شود؛ وضعیتی که در آن عامل یاد می‌گیرد برای دریافت پاداش کمکی بیشتر، کارهای تکراری و بی‌فایده انجام دهد بدون آنکه به هدف اصلی محیط برسد. آن‌ها اثبات کردند که برای تضمین ثبات و حفظ سیاست بهینه محیط (\lr{Policy Invariance})، پاداش کمکی $F$ باید حتماً به صورت تفاضل پتانسیل یک تابع حالت پایه‌ریزی شود. این متدولوژی تحت عنوان شکل‌دهی پاداش مبتنی بر پتانسیل (\lr{Potential-Based Reward Shaping - PBRS}) شناخته می‌شود \cite{ng1999policy} و فرمولاسیون ریاضی آن به شرح زیر است:
\begin{equation}
    F(s, a, s') = \gamma \Phi(s') - \Phi(s)
\end{equation}
که در آن $\gamma$ ضریب تخفیف مارکوف و $\Phi(s)$ تابع پتانسیل حالت $s$ است. بر اساس قضیه ریاضی نگ، هرگاه پاداش شکل‌دهی دارای این فرمولاسیون باشد، مجموعه سیاست‌های بهینه عامل $\pi^*$ تحت پاداش شکل‌دهی‌شده دقیقاً با مجموعه سیاست‌های بهینه تحت پاداش اصلی یکسان خواهد بود\footnote{قضیه نگ \cite{ng1999policy} شرط کافی برای حفظ سیاست بهینه را ارائه می‌دهد، اما فرض می‌کند تابع پتانسیل $\Phi$ ایستا (\lr{static}) است. در سیستم‌های خودوفق با رانش مفهومی، این فرض نقض می‌شود. مکانیسم شکل‌دهی تصادفی \lr{RS-DRL} (فصل ۴، بخش ۴-۶) نیازی به تابع پتانسیل ایستا ندارد و بنابراین از این محدودیت تئوریک عبور می‌کند.}.

\textbf{نقد و بررسی انتقادی:}
اگرچه \lr{PBRS} چارچوب ریاضی مستحکمی برای تضمین حفظ سیاست بهینه ارائه می‌دهد، اما استقرار آن در سیستم‌های خودوفق واقعی با مانع بزرگی مواجه است.
محدودیت نخست، \textbf{استاتیک بودن تابع پتانسیل در محیط‌های پویا} است؛ در طراحی سنتی \lr{PBRS} \cite{ng1999policy}، تابع پتانسیل $\Phi(s)$ به صورت کاملاً ایستا (\lr{Static}) در زمان طراحی سیستم بر اساس دانش تجربی کارشناسان دامنه تعریف می‌شود.
محدودیت دوم، \textbf{عدم انطباق با رانش مفهومی و شرایط پویای آنلاین} است؛ در یک سیستم خودوفق تحت تأثیر عدم قطعیت مستمر (رانش مفهومی)، پتانسیل واقعی حالت‌ها دائماً تغییر می‌کند. یک پتانسیل ایستای تعریف‌شده در زمان طراحی، قادر به بازتاب پویایی‌های پیش‌بینی‌نشده زمان اجرا نیست و در شرایط جدید، پاداش‌های کمکی گمراه‌کننده‌ای تولید می‌کند که سرعت یادگیری عامل را به شدت کاهش داده یا همگرایی آنلاین را مختل می‌سازد \cite{gheibi2024dealing}.

این نقیصه بنیادین تئوریک، ضرورت ابداع یک فرآیند شکل‌دهی تصادفی و پویای پاداش (همچون الگوریتم پیشنهادی \lr{RS-DRL} در این رساله) را آشکار می‌سازد تا پتانسیل حالت‌ها را به صورت زمان اجرا و آنلاین بر اساس تضمین‌های کیفی محاسبات صوری تنظیم نماید.

\section{ارزیابی و تایید صوری در فرآیند خودوفق‌سازی}
\label{sec:ch3_formal_verification}

تضمین ایمنی و ارضای کیفیت خدمات (\lr{QoS}) در سیستم‌های خودوفق، به ویژه در حوزه‌های حساس، نیازمند ابزارهای صوری مستحکم است. ارزیابی و تأیید صوری زمان اجرا (\lr{Runtime Formal Verification}) به عنوان پارادایمی مطرح است که به سیستم اجازه می‌دهد تا به صورت ریاضیاتی و دقیق، صحت تصمیمات خود را پیش از اعمال به سیستم واقعی ارزیابی کند. در این بخش، مبانی نظری تأیید صوری کمّی زمان اجرا، ابزارها و چارچوب‌های شاخص (مانند \lr{ActivFORMS} و \lr{UPPAAL-SMC}) و در نهایت معماری‌های ترکیبی (ترادف) صوری-یادگیری را مورد تحلیل انتقادی قرار می‌دهیم.

\subsection{تأیید صوری کمّی زمان اجرا \lr{(QVR)}}
\label{sec:ch3_qvr}

تأیید صوری کمّی زمان اجرا (\lr{Quantitative Verification at Runtime - QVR}) که توسط پیشگامانی نظیر کالینسکو (\lr{Calinescu}) و همکارانش \cite{calinescu2011formal} و وینز (\lr{Weyns}) \cite{weyns2012formal} توسعه یافت، فرآیندی است که در آن مشخصات غیرعملکردی سیستم (مانند کارایی، استواری و مصرف انرژی) در زمان اجرا با استفاده از مدل‌های ریاضی احتمالات ارزیابی می‌شوند.

در این پارادایم، پویایی‌های سیستم و عدم قطعیت‌های محیطی معمولاً به صورت زنجیره‌های مارکوف زمان‌گسسته (\lr{DTMC}) یا فرآیندهای تصمیم‌گیری مارکوف (\lr{MDP}) مدل‌سازی می‌شوند. اهداف کیفی سیستم نیز با استفاده از منطق‌های زمانی احتمالات نظیر منطق زمانی محاسبات احتمالات (\lr{PCTL}) فرموله می‌گردند. برای مثال، هدف ایمنی شبکه به صورت فرمول زیر تعریف می‌شود:
\begin{equation}
    P_{\geq 0.95} [ F^{\leq t} (\text{PacketLoss} < 0.1) ]
\end{equation}
که به معنای آن است که با احتمال حداقل ۹۵ درصد، نرخ گم‌شدن بسته‌ها در طول زمان $t$ کمتر از ۱۰ درصد باقی بماند. موتورهای ارزیابی صوری زمان اجرا با حل ریاضیاتی این مدل‌ها به صورت برخط، مشخص می‌سازند که آیا پیکربندی انتخاب‌شده قیود فوق را ارضا می‌کند یا خیر \cite{calinescu2011formal}.

\subsection{چارچوب‌های ActivFORMS و ابزار UPPAAL-SMC}
\label{sec:ch3_activforms_uppaal}

در بین چارچوب‌های صوری توسعه‌یافته برای سیستم‌های خودوفق، چارچوب \lr{ActivFORMS} (\lr{Active Formal Models for Self-Adaptation}) که توسط وینز و همکارانش \cite{weyns2018activforms} معرفی شد، جایگاه ویژه‌ای دارد. \lr{ActivFORMS} یک موتور مبتنی بر مدل‌های صوری اجرایی در زمان اجراست که درون حلقه \lr{MAPE-K} قرار می‌گیرد. این چارچوب تضمین می‌کند که فرآیند وفقی همواره با مدل‌های صوری اعتبارسنجی‌شده مطابقت داشته و رفتارهای غیرایمن در سیستم رخ ندهند.

با این حال، حل کامل و جامع توابع احتمالاتی در زمان اجرا به شدت پرهزینه است. برای عبور از این محدودیت محاسباتی، استفاده از ابزار بررسی مدل آماری (\lr{Statistical Model Checking - SMC}) تحت چارچوب \lr{UPPAAL-SMC} \cite{david2015uppaal} مورد توجه قرار گرفت. الگوریتم‌های \lr{SMC} به جای حل فرمول‌های ریاضی پیچیده در تمام فضای حالت (که منجر به انفجار فضا می‌شود)، از شبیه‌سازی‌های تصادفی هوشمند سیستم بهره می‌برند. \lr{UPPAAL-SMC} با اجرای تعداد مشخصی شبیه‌سازی زمان اجرا و اعمال آزمون‌های فرض آماری صوری (مانند آزمون نسبت احتمال متوالی والد - \lr{Wald's Sequential Probability Ratio Test - SPRT})، ارضای یا عدم ارضای خواص را با یک سطح اطمینان آماری مشخص ($\alpha, \beta$) تعیین می‌کند. این مکانیزم به شدت سربار محاسباتی ارزیابی‌های صوری زمان اجرا را کاهش می‌دهد.

\subsection{معماری‌های ترکیبی یادگیری و ارزیابی صوری (مدل‌های ترادف)}
\label{sec:ch3_hybrid_verification_learning}

با آشکار شدن قدرت مقیاس‌پذیری یادگیری ماشین و امنیت بالای روش‌های صوری، محققان در سال‌های اخیر به سمت توسعه معماری‌های ترکیبی متمایل شده‌اند که در ادبیات علمی به عنوان رویکردهای ترادف (\lr{Tandem Approaches}) شناخته می‌شوند. پژوهش کامارا و همکارانش \cite{camara2020quantitative} نمونه شاخصی از این معماری‌هاست. در این رویکردها، عامل یادگیری تقویتی وظیفه هدایت سیستم و یادگیری سیاست‌های وفقی را بر عهده دارد، در حالی که ماژول ارزیابی صوری زمان اجرا (\lr{QVR}/\lr{SMC}) به عنوان یک فیلتر ایمنی زمان اجرا عمل می‌کند. هرگاه اقدام پیشنهادی عامل یادگیرنده از نظر صوری ناامن تشخیص داده شود، فیلتر صوری آن را بلاک کرده و سیستم را مجبور به انتخاب گزینه‌ای ایمن می‌نماید.

\textbf{نقد و بررسی انتقادی:}
اگرچه معماری‌های ترادفی کلاسیک گام بزرگی رو به جلو هستند، اما کماکان دارای یک نقطه ضعف عملیاتی جدی هستند که مانع از استقرار آن‌ها در بسترهای اینترنت اشیا بزرگ می‌گردد.
نقطه ضعف اول، \textbf{سربار شدید ارزیابی صوری در فضاهای وفقی بزرگ} است؛ بررسی مدل صوری (حتی با ابزار سبک \lr{UPPAAL-SMC}) برای هر یک از گزینه‌های تصمیم‌گیری زمان اجرا مستلزم چند صد بار شبیه‌سازی تصادفی است. در صورتی که عامل یادگیرنده بخواهر در یک فضای بزرگ (مثلاً با ۴۰۹۶ گزینه وفقی) کاوش کند، ارزیابی صوری تمامی این گزینه‌ها زمان تصمیم‌گیری حلقه \lr{MAPE-K} را به شدت افزایش داده (تا چندین دقیقه) و عملاً کارکرد آنلاین سیستم را با اختلال جدی مواجه می‌سازد.
نقطه ضعف دوم، \textbf{ایستا بودن فیلتر ایمنی صوری} است؛ فیلترهای صوری کلاسیک رفتاری صفر و یکی دارند و نمی‌توانند بازخورد کیفی دقیقی به عامل یادگیرنده ارائه دهند تا او سرعت یادگیری خود را بهبود بخشد.

این نقد اساسی، ضرورت وجود یک معماری جدید مانند \lr{RS-DRL} را اثبات می‌کند\footnote{معماری \lr{RS-DRL} با جداسازی زمانی (\lr{Temporal Separation}) که در فصل ۴ (بخش ۴-۲ و ۴-۴) تشریح شده است، این مشکل را حل می‌کند: ارزیابی صوری به صورت موازی و ناهمگام (\lr{asynchronous}) در پس‌زمینه انجام می‌شود و مسیر تصمیم‌گیری آنلاین را مسدود نمی‌کند. شکل ۴-۷ در فصل ۴ این معماری غیرانسدادی را نشان می‌دهد.}؛ معماری خاصی که در آن ماژول صوری \lr{UPPAAL-SMC} به جای ایفای نقش به عنوان فیلتر نهایی سنگین برای تمامی گزینه‌ها، به صورت ناهمگام (\lr{Asynchronous}) برای محاسبه‌ی صوری متغیرهای پتانسیل یادگیری و شکل‌دهی تصادفی پاداش به کار گرفته شود. این کار موازنه بهینه‌ای بین مقیاس‌پذیری تصمیم‌گیری عمیق و تضمین‌های ریاضی روش‌های صوری برقرار می‌سازد.

\section{سنتز: شکاف تحقیقاتی سه‌گانه و موقعیت‌دهی علمی RS-DRL}
\label{sec:ch3_synthesis}

بر اساس نقد و بررسی‌های انجام‌شده در بخش‌های پیشین، مشخص می‌گردد که هر یک از رویکردهای کلاسیک و مدرن مدیریت فضای وفقی در سیستم‌های خودوفق، از مزایای متمایزی برخوردار بوده و همزمان با محدودیت‌های بنیادین تئوریک و عملی مواجه هستند. به منظور جمع‌بندی دقیق وضعیت هنر، در ادامه ابتدا یک ماتریس ارزیابی چندبعدی ارائه گردیده و سپس شکاف‌های تحقیقاتی موجود تبیین می‌شوند.

\subsection{ماتریس ارزیابی و سنتز پیشینه پژوهش}
\label{sec:ch3_synthesis_matrix}

در جدول \ref{tab:ch3_synthesis}، پارادایم‌های اصلی موجود در پیشینه پژوهش بر اساس چارچوب ارزیابی پنج‌بعدی پیشنهادی به همراه روش پیشنهادی این رساله (\lr{RS-DRL}) مورد مقایسه و ارزیابی تطبیقی قرار گرفته‌اند.

\begin{table*}[htbp]
\centering
\caption{مقایسه تطبیقی پارادایم‌های مدیریت فضای وفقی زمان اجرا بر اساس چارچوب ارزیابی پنج‌بعدی (برگرفته از تحلیل کارهای پیشین نظیر \cite{coker2015sass, kinneer2021information, du2023scalability, gheibi2024dealing})}
\label{tab:ch3_synthesis}
\resizebox{\textwidth}{!}{%
\begin{tabular}{|p{2.6cm}|p{2.3cm}|p{2.3cm}|p{2.3cm}|p{2.3cm}|p{2.3cm}|}
\hline
\textbf{پارادایم / رویکرد} & \textbf{مدیریت ابعاد بزرگ (مقیاس‌پذیری)} & \textbf{استواری در برابر تغییرات آنلاین} & \textbf{استقلال از داده و پیش‌آموزش} & \textbf{کارایی محاسباتی آنلاین} & \textbf{فرار از تله بهینه‌های محلی} \\ \hline
\textbf{هرس مبتنی بر یادگیری ماشین} \cite{quin2019efficient, kachi2023hybrid} & متوسط & ضعیف (به علت رانش مفهومی) & ضعیف (نیازمند برچسب‌گذاری آفلاین) & عالی (میلی‌ثانیه) & غیرقابل اعمال (فقط هرس می‌کند) \\ \hline
\textbf{جستجوی آنلاین تکاملی} \cite{coker2015sass, kinneer2020building, chen2018femosaa} & ضعیف (انفجار زمان محاسبه) & متوسط & عالی (بدون نیاز به پیش‌آموزش) & ضعیف (ثانیه/دقیقه) & متوسط (به واسطه جهش‌های تصادفی) \\ \hline
\textbf{یادگیری تقویتی عمیق سنتی} \cite{kim2009reinforcement, du2023scalability} & متوسط تا عالی & متوسط & عالی (یادگیری آنلاین مدل‌فردا) & عالی (کمتر از چند میلی‌ثانیه) & ضعیف (به علت کاوش کورکورانه \lr{$\epsilon$-greedy}) \\ \hline
\textbf{تأیید صوری زمان اجرا} \cite{weyns2018activforms, camara2020quantitative} & بسیار ضعیف (انفجار فضا) & عالی (تضمین ریاضی پایداری) & عالی & بسیار ضعیف (سربار شدید محاسبات صوری) & متوسط (تنها ایمنی را بررسی می‌کند) \\ \hline
\textbf{روش پیشنهادی (RS-DRL)} & \textbf{عالی (استفاده از شبکه عمیق بدون فیلتر برخط)} & \textbf{عالی (انطباق با تغییرات به واسطه پتانسیل پویا)} & \textbf{عالی (یادگیری کاملاً آنلاین و مدل‌فردا)} & \textbf{عالی (کمتر از چند میلی‌ثانیه برای استنتاج)} & \textbf{عالی (به واسطه شکل‌دهی پویای پاداش صوری)} \\ \hline
\end{tabular}%
}
\end{table*}

\subsection{شکاف تحقیقاتی سه‌گانه}
\label{sec:ch3_triple_gap}

با متقاطع‌سازی ابعاد ارزیابی و چالش‌های عملیاتی احصاءشده از پیشینه پژوهش، \textbf{شکاف تحقیقاتی سه‌گانه‌ای} که این رساله بر حل آن‌ها تمرکز دارد به شرح زیر صورت‌بندی می‌شود:

\begin{enumerate}
    \item \textbf{شکاف اول: تقابل مقیاس‌پذیری و پایداری زمانی تصمیم‌گیری (\lr{Scale-Latency Dilemma}):}
    در سیستم‌های خودوفق با فضاهای وفقی بزرگ (بیش از ۱۰۰۰ گزینه)، رویکردهایی که از ارزیابی‌های صوری دقیق زمان اجرا (\lr{QVR}/\lr{SMC}) یا جستجوهای تکاملی استفاده می‌کنند، با چالش جدی انفجار محاسباتی مواجه هستند که زمان تصمیم‌گیری را به مقیاس دقیقه افزایش می‌دهد. در مقابل، روش‌های یادگیری ماشین و شبکه‌های عمیق زمان استنتاج میلی‌ثانیه‌ای دارند اما فاقد هرگونه تضمین ریاضی ثبات و ایمنی آنلاین هستند. تا کنون مکانیزمی خلاقانه که بتواند بدون تحمیل سربار محاسباتی صوری به مسیر استنتاج آنلاین، تضمین‌های صوری ریاضی را به یادگیری عمیق تزریق کند، ارائه نشده است.
    
    \item \textbf{شکاف دوم: عدم تطابق مکانیزم‌های نگاه به گذشته با اهداف آستانه‌محور خودوفقی (\lr{Goal-Threshold Mismatch}):}
    الگوریتم‌های پیشرفته یادگیری برای پاداش‌های پراکنده (مانند \lr{HER} و \lr{Soft HER}) منحصراً برای محیط‌های «هدف‌محور» طراحی شده‌اند که عامل در آن‌ها به دنبال رسیدن به مختصاتی معین است. این در حالی است که اهداف سیستم‌های خودوفق معاصر ماهیت «آستانه‌محور» (مانند نگه‌داشتن مصرف انرژی زیر یک مقدار مشخص) دارند. این ناسازگاری تئوریک باعث می‌شود که در صورت بروز نقض آستانه، برچسب‌گذاری مجدد با نگاه به گذشته فاقد اعتبار ریاضی گشته و عملاً کارایی خود را در هدایت کاوش آنلاین سیستم‌های خودوفق از دست بدهد.
    
    \item \textbf{شکاف سوم: ایستایی توابع پتانسیل در مواجهه با پویایی‌های نویزی آنلاین (\lr{Static-Dynamic Potential Gap}):}
    تکنیک‌های ریاضیاتی شکل‌دهی پاداش مبتنی بر پتانسیل (\lr{PBRS}) \cite{ng1999policy} که ثبات سیاست بهینه را تضمین می‌کنند، به صورت سنتی نیازمند یک تابع پتانسیل $\Phi(s)$ کاملاً ایستا هستند که در زمان طراحی توسط کارشناس خبره تعریف می‌شود. در حضور عدم قطعیت‌های شدید زمان اجرا و رانش مفهومی محیط، این توابع ایستای زمان طراحی کارایی خود را از دست داده و حتی ممکن است با ارائه پاداش‌های نادرست، همگرایی آنلاین عامل را مختل سازند. خلأ شدیدی در زمینه سیستم‌های شکل‌دهی پویای پتانسیل وجود دارد که بتوانند پتانسیل حالت‌ها را به صورت پویا بر اساس بازخوردهای صوری آنلاین به‌روزرسانی کنند.
\end{enumerate}

\subsection{موقعیت‌دهی علمی معماری پیشنهادی RS-DRL}
\label{sec:ch3_rsdrl_positioning}

معماری پیشنهادی این رساله، تحت عنوان \textbf{یادگیری تقویتی عمیق شکل‌دهی‌شده با بررسی صوری زمان اجرا (\lr{RS-DRL})}، به عنوان راهکاری جامع برای پل زدن بر این شکاف‌های سه‌گانه طراحی و موقعیت‌دهی شده است. نوآوری‌های کلیدی این رویکرد به شرح زیر با شکاف‌های تحقیقاتی متناظر همخوانی دارند:

\begin{itemize}
    \item \textbf{پاسخ به شکاف اول (تقابل سرعت و امنیت):} معماری \lr{RS-DRL} فرآیند زمان‌بر ارزیابی صوری زمان اجرا را از مسیر مستقیم تصمیم‌گیری آنلاین عامل یادگیرنده خارج ساخته و آن را به صورت \textbf{غیرهمگام (\lr{Asynchronous})} در پس‌زمینه اجرا می‌کند. عامل برای تصمیم‌گیری صرفاً شبکه عمیق را در کمتر از ۱ میلی‌ثانیه استنتاج می‌کند، در حالی که ماژول صوری تصادفی (\lr{UPPAAL-SMC}) به صورت موازی پتانسیل‌های یادگیری را به‌روزرسانی می‌کند. این متدولوژی برای اولین بار به طور همزمان مقیاس‌پذیری تصمیم‌گیری عمیق آنلاین و استواری صوری ریاضیاتی را تضمین می‌نماید.

    \item \textbf{پاسخ به شکاف دوم (یادگیری آستانه‌محور):} در \lr{RS-DRL}، به جای به‌کارگیری متدهای هدف‌محور ناکارآمد مانند \lr{HER}، سیستم با ترکیب پاداش محیط و پاداش کمکی صوری هدایت می‌شود. این پاداش کمکی به طور دقیق شدت ارضای آستانه‌ها را منعکس کرده و بدون نیاز به تعریف مختصات هدف صریح، امنیت کاوش عامل را در فضاهای آستانه‌محور تضمین می‌نماید.

    \item \textbf{پاسخ به شکاف سوم (شکل‌دهی پویا):} نوآوری بنیادین ریاضیاتی در \lr{RS-DRL}، معرفی استراتژی \textbf{شکل‌دهی پویای پاداش مبتنی بر بررسی صوری زمان اجرا} است. در این متدولوژی، تابع پتانسیل دیگر ایستا نیست، بلکه به عنوان یک متغیر پویای تصادفی آنلاین فرموله‌بندی شده و مقادیر آن به طور مستمر بر اساس خروجی‌های کمّی حاصل از تحلیل‌های آماری \lr{UPPAAL-SMC} به‌روزرسانی می‌گردند. این مکانیزم انطباق کامل عامل با رانش مفهومی و نویز محیط را بدون ابطال تئوریک سیاست بهینه رقم می‌زند.
\end{itemize}

جدول \ref{tab:ch3_positioning} خلاصه‌ای از نحوه پوشش شکاف‌های تحقیقاتی توسط نوآوری‌های \lr{RS-DRL} را نمایش می‌دهد.

\begin{table}[htbp]
\centering
\caption{نحوه پوشش شکاف‌های تحقیقاتی توسط نوآوری‌های معماری پیشنهادی RS-DRL}
\label{tab:ch3_positioning}
\resizebox{\textwidth}{!}{%
\begin{tabular}{|p{3cm}|p{5cm}|p{6cm}|}
\hline
\textbf{شکاف تحقیقاتی} & \textbf{محدودیت روش‌های کلاسیک} & \textbf{راهکار و نوآوری ارائه شده در RS-DRL} \\ \hline
\textbf{شکاف اول (سرعت و امنیت)} & تأخیر طولانی مدل‌چکینگ زمان اجرا در فضاهای بزرگ & ارزیابی صوری ناهمگام و استنتاج میلی‌ثانیه‌ای آنلاین شبکه عمیق \\ \hline
\textbf{شکاف دوم (یادگیری آستانه‌محور)} & عدم امکان استفاده از روش‌های برچسب‌گذاری مجدد هدف \lr{(HER)} & هدایت عامل به کمک پاداش‌های کمکی صوری سنجش کیفیت و آستانه \\ \hline
\textbf{شکاف سوم (شکل‌دهی پویا)} & ایستایی توابع پتانسیل شکل‌دهی پاداش در محیط‌های متغیر آنلاین & شکل‌دهی پویای پاداش با به‌روزرسانی آنلاین پتانسیل‌ها به وسیله بازخورد صوری زمان اجرا \\ \hline
\end{tabular}%
}
\end{table}

\section{خلاصه فصل}
\label{sec:ch3_summary}

در این فصل، مرور عمیق و منسجمی بر مبانی نظری و پیشینه پژوهش مرتبط با فرآیندهای خودوفق‌سازی در سیستم‌های نرم‌افزاری و شبکه‌های توزیع‌شده نوین صورت گرفت. مسیر تحلیل پیشینه با تبیین ماهیت معماری‌های خودوفق مبتنی بر حلقه \lr{MAPE-K} و چالش بنیادین انفجار فضای پیکربندی در سیستم‌های مقیاس‌بزرگ آغاز گردید. سپس، رویکردهای موجود در وضعیت هنر در قالب چهار پارادایم اصلی طبقه‌بندی و مورد بررسی انتقادی قرار گرفتند: هرس مبتنی بر یادگیری ماشین سنتی، جستجوی آنلاین تکاملی، یادگیری تقویتی عمیق بدون مدل و تأیید صوری زمان اجرا.

با اعمال یک چارچوب ارزیابی پنج‌بعدی پیشنهادی بر پیشینه پژوهش، مشخص گردید که هیچ‌یک از متدولوژی‌های کلاسیک قادر به حل همزمان چالش‌های مقیاس‌پذیری، کارایی محاسباتی آنلاین، استواری در برابر رانش مفهومی و تضمین‌های ریاضیاتی ایمنی نیستند. این تحلیل انتقادی به صورت‌بندی دقیق \textbf{شکاف تحقیقاتی سه‌گانه‌ای} منجر شد که محور تمرکز علمی این رساله را تشکیل می‌دهد:
\begin{enumerate}
    \item \textbf{تقابل مقیاس‌پذیری و پایداری زمانی تصمیم‌گیری:} ناشی از سربار زمانی شدید مدل‌چکینگ آنلاین صوری در فضاهای پیکربندی بزرگ.
    \item \textbf{عدم تطابق مکانیزم‌های نگاه به گذشته با اهداف آستانه‌محور خودوفقی:} به دلیل محدودیت ذاتی الگوریتم‌هایی مانند \lr{HER} که منحصراً برای فضاهای وفقی مختصات‌محور و غایت‌محور توسعه یافته‌اند.
    \item \textbf{ایستایی توابع پتانسیل شکل‌دهی پاداش:} که با ماهیت نویزی و دائماً در حال تغییر عدم قطعیت‌های آنلاین زمان اجرا ناسازگار است.
\end{enumerate}

در نهایت، موقعیت‌دهی علمی روش پیشنهادی این رساله (\lr{RS-DRL}) به عنوان پاسخ و راهکاری یکپارچه برای پوشش هر سه شکاف فوق تبیین شد. بنابراین، یافته‌های تحلیلی این فصل به طور قاطع تأیید می‌کنند که هیچ‌یک از رویکردهای موجود در پیشینه پژوهش قادر به ارضای همزمان الزامات سه‌گانه (سرعت میلی‌ثانیه‌ای، انطباق‌پذیری با پویایی زمان اجرا، و تضمین صوری ایمنی ریاضی) نیستند و معماری پیشنهادی \lr{RS-DRL} دقیقاً برای پر کردن این شکاف بنیادین علمی موقعیت‌دهی شده است.

\begin{framed}
\noindent\textbf{خلاصه یافته‌های کلیدی فصل سوم (مرور پیشینه پژوهش):}
\begin{itemize}
    \item \textbf{ناکارآمدی فیلترهای صوری آنلاین:} فیلترهای صوری زمان اجرا (\lr{QVR}) به صورت مستقیم در مسیر تصمیم‌گیری باعث سربار محاسباتی شدید (در حد دقیقه) برای فضاهای وفقی بزرگ می‌گردند \cite{camara2020quantitative, calinescu2011formal}.
    \item \textbf{انفجار فضای حالت در تکامل آنلاین:} روش‌های جستجوی تصادفی و تکاملی آنلاین فاقد توانایی تصمیم‌گیری بی‌درنگ در مقیاس میلی‌ثانیه برای سیستم‌های اینترنت اشیاء هستند \cite{coker2015sass, kinneer2021information}.
    \item \textbf{شکست روش‌های HER در سناریوهای SLA:} الگوریتم‌های برچسب‌گذاری مجدد هدف به علت تعاریف مبتنی بر فیزیک مختصاتی، در برابر آستانه‌های کیفی سیستم‌های خودوفق فاقد کارایی علمی می‌باشند.
    \item \textbf{ضرورت معماری ترادف ناهمگام:} تلفیق ناهمگام یادگیری عمیق و بررسی صوری آماری (\lr{UPPAAL-SMC}) به عنوان تنها راهکار تأمین همزمان کارایی آنلاین و تضمین‌های ریاضیاتی ایمنی شناسایی شد.
\end{itemize}
\end{framed}

\textbf{پل مفهومی به فصل بعد (بیانیه پل):}
با استقرار بستر نظری مستحکم در این فصل و صورت‌بندی دقیق نقاط ضعف و شکاف‌های پیشینه، مسیر برای تشریح جزئیات فنی و عملیاتی متدولوژی پیشنهادی هموار گردیده است. بر این اساس، در \textbf{فصل چهارم} به تبیین دقیق معماری و سازماندهی فرآیند اجرایی الگوریتم \lr{RS-DRL} خواهیم پرداخت. در آن فصل، فرمول‌بندی ریاضیاتی شکل‌دهی پویای پاداش صوری، نحوه ارتباط ناهمگام میان عامل یادگیری تقویتی عمیق و ماژول بررسی صوری احتمالات (\lr{UPPAAL-SMC}) و در نهایت، تضمین همگرایی تئوریک سیستم تحت شرایط نویزی زمان اجرا تشریح خواهند شد تا نشان داده شود چگونه معماری پیشنهادی قادر است موازنه بهینه‌ای بین سرعت میلی‌ثانیه‌ای آنلاین و ایمنی صوری ریاضیاتی برقرار سازد.
\newpage


% ==================== chapters/chapter4.tex ====================


\bibliographystyle{ieeetr-fa}
\bibliography{References,MyReferences}
\end{document}
